% !TEX root = ../main.tex
% !TEX program = xelatex
\title{基于多源预测与线性规划的微网协同调度}
\date{}
\maketitle
\begin{abstract}
\noindent\hspace*{2em}针对含光伏、储能和外部电网的小区微网购电调度问题，本文以10分钟为基本时段，
建立统一的电量流与储能状态模型，并依次研究确定性调度、供需不确定、日内预报更新和波动电价下的购电策略。

\textbf{针对问题一}，建立\textbf{全时段电量流线性规划模型}，统筹外网购电、光伏利用和储能充放电。
代表日购电费用为 35126.85元，较峰谷策略和光伏空间预留策略分别降低 \textbf{8.23\%} 和 \textbf{7.43\%}。

\textbf{针对问题二}，建立\textbf{周期预测—近期加权准备—跨日线性规划模型}。
分别预测负载和光伏，并根据近期净需求误差设置安全准备量，以降低5倍紧急购电风险。
2025年2月至12月总费用为 \textbf{1358.38万元}，较等权准备方案节省 44.66万元。

\textbf{针对问题三}，建立\textbf{融合预测与日内校正滚动规划}。
用近期同星期数据预测负载，再根据当天已发生的偏差修正后续判断，并融合官方光伏预报。
结合安全准备与实际库存，滚动调整未执行订单，总费用为 \textbf{1301.42万元}。
同一模型下的三组信息对照表明，日内反馈与改单节省50.63万元，后发官方预报进一步节省 \textbf{6.95万元}。

\textbf{针对问题四}，将波动电价纳入预测与调度。
零点一次订电采用\textbf{提升树预测与零点规划}，日内调整采用\textbf{岭回归预测与日内滚动规划}，
两种方案总费用分别为 \textbf{1424.86万元} 和 \textbf{1369.59万元}。

模型检验表明，各方案均满足供需平衡及储能容量、功率约束。
本文形成了由\textbf{确定性调度—供需不确定性—日内信息更新—价格不确定性}逐步扩展的微网购电优化框架，可为类似微网的滚动购电与储能调度提供参考。
\end{abstract}
\noindent{\heiti\zihao{-4} 关键词：}{\songti\zihao{-4} 微网调度；储能优化；线性规划；滚动优化；波动电价}
